% Part: proof-theory
% Chapter: cut-elimination
% Section: midsequent

\documentclass[../../../include/open-logic-section]{subfiles}

\begin{document}

\olfileid{pt}{cut}{mid}

\olsection{মধ্যবর্তী সিকোয়েন্ট উপপাদ্য ও Herbrand-এর উপপাদ্য}

কর্তন-বর্জন উপপাদ্যের একটি গুরুত্বপূর্ণ ফল হলো মধ্যবর্তী
সিকোয়েন্ট উপপাদ্য। এটি বলে: $\Gamma \Sequent \Delta$-র
এমন !!a{proof}~$\pi$ থাকলে, যাতে প্রতিটি !!{formula}
অগ্র-পরিমাণসূচক রূপে আছে, এমন !!a{proof}~$\pi'$-ও আছে যাতে
$S \ident \Pi \Sequent \Lambda$ নামে একটি সিকোয়েন্ট
(\emph{মধ্যবর্তী সিকোয়েন্ট}) থাকে; $\pi'$-তে~$S$-এর উপরের
সব অনুমান বচনীয়, আর $S$-এর নীচের সব অনুমান পরিমাণসূচক-সংক্রান্ত।
(!!^a{formula}~$!A$ \emph{অগ্র-পরিমাণসূচক} বা
\emph{অগ্র-পরিমাণসূচক রূপে} থাকে যদি তার রূপ
\[\mathsf{Q}_1x_1\dots\mathsf{Q}_nx_n\,!B(x_1,\dots,x_n)\] হয়,
যেখানে প্রত্যেক $\mathsf{Q}_i$ হয় $\lforall$, নয়তো~$\lexists$।)
মধ্যবর্তী সিকোয়েন্ট উপপাদ্য থেকে Herbrand-এর উপপাদ্যও
অনুসৃত হয়। এর সাধারণ ক্ষেত্রটি বর্ণনা করা কিছুটা কঠিন;
মোটামুটি ফলটির রূপ এমন: প্রথম-ক্রম যুক্তির প্রতিটি
!!{provable} !!{sentence}~$!A$-র জন্য এমন একটি বচনীয়
সর্বতঃসত্য সূত্র~$!A'$ আছে, যা থেকে $!A$ !!{prove} করা যায়।
এখানে $\lexists[x][!B(x)]$ রূপের !!{sentence}s-এর
ক্ষেত্রটি দিই, যেখানে $!B$ পরিমাণসূচকবিহীন:

\begin{thm}[Herbrand-এর উপপাদ্য]
ধরি, $!B(x)$ পরিমাণসূচকবিহীন এবং $\Proves
\lexists[x][!B(x)]$। তাহলে এমন পদ~$t_1$, \dots, $t_n$
আছে যাতে $\Proves !B(t_1) \lor \dots \lor !B(t_n)$।
\end{thm}

!!{formula}~$!B(t_1) \lor \dots \lor !B(t_n)$-কে
$\lexists[x][!B(x)]$-এর \emph{Herbrand-বিচ্ছেদ} বলে।
$!B$ পরিমাণসূচকবিহীন বলে এই বিচ্ছেদও তাই। সুতরাং এটি
!!{provable} হলে কর্তনমুক্ত প্রমাণে !!{provable}, ফলে
শুধু বচনীয় অনুমান-ব্যবহারকারী !!a{proof}-ও আছে।
অতএব এটি সর্বতঃসত্য।

Herbrand-এর উপপাদ্যের কঠিন অংশ হলো, প্রতিটি !!{provable}
!!{formula}~$\lexists[x][!B(x)]$-এর জন্য একটি
Herbrand-বিচ্ছেদের অস্তিত্ব দেখানো। বিপরীতটি সহজ:
$\lexists[x][!B(x)]$-এর Herbrand-বিচ্ছেদ থাকলে সূত্রটি
প্রমাণযোগ্য। আসলে, বিচ্ছেদটি থেকে $\lexists[x][!B(x)]$
অনুসৃত হয়। যেমন, $n=2$ হলে $\Log{G3c}$-তে এই !!a{proof}:
\[
\Axiom$!B(t_1) \fCenter \lexists[x][!B(x)], !B(t_1), !B(t_2)$
\Axiom$!B(t_2) \fCenter \lexists[x][!B(x)], !B(t_1), !B(t_2)$
\RightLabel{\LeftR{\lor}}
\BinaryInf$!B(t_1) \lor !B(t_2) \fCenter \lexists[x][!B(x)], !B(t_1), !B(t_2)$
\RightLabel{\RightR{\lexists}}
\UnaryInf$!B(t_1) \lor !B(t_2) \fCenter \lexists[x][!B(x)], !B(t_2)$
\RightLabel{\RightR{\lexists}}
\UnaryInf$!B(t_1) \lor !B(t_2) \fCenter \lexists[x][!B(x)]$
\DisplayProof
\]

$\Sequent \lexists[x][!B(x)]$-র !!a{proof}~$\pi$
থেকে Herbrand-বিচ্ছেদ বের করতে প্রথমে এমন ক্ষেত্র ধরি,
যেখানে কর্তনমুক্ত !!{proof}~$\pi$-র সব
$\RightR{\lexists}$ অনুমান শেষে আছে:
\[
\AxiomC{}
\RightLabel{$\pi_1$}
\Deduce$ \fCenter \lexists[x][!B(x)], !B(t_1), \dots, !B(t_n)$
\RightLabel{\RightR{\lexists}}
\UnaryInf$ \fCenter \lexists[x][!B(x)], !B(t_2), \dots, !B(t_n)$
\RightLabel{$(n-1) \times \RightR{\lexists}$}
\Deduce$ \fCenter \lexists[x][!B(x)]$
\DisplayProof
\]
$\pi$-তে $\lexists[x][!B(x)]$ সক্রিয় এমন সব
$\RightR{\lexists}$ অনুমান এগুলিই হলে উপপ্রমাণ~$\pi_1$-তে
আর কোনো পরিমাণসূচক অনুমান নেই: $!B(x)$-তে অন্য
পরিমাণসূচক নেই এবং অন্তিম সিকোয়েন্টের বাম পাশে কোনো
!!{formula}s নেই। অন্য কথায়,
$ \fCenter \lexists[x][!B(x)], !B(t_1), \dots, !B(t_n)$
হলো~$\pi$-র মধ্যবর্তী সিকোয়েন্ট। তার উপরে পরিমাণসূচক
অনুমান না থাকায় উপরের সব সিকোয়েন্টে
$\lexists[x][!B(x)]$ থাকলেও তা সক্রিয় বা প্রধান নয়।
তাই !!{proof}~$\pi_1$ থেকে ওই সূত্র বাদ দিয়ে
$\fCenter !B(t_1), \dots, !B(t_n)$-র !!a{proof} পাই;
তারপর $\RightR{\lor}$-এর $n-1$টি প্রয়োগে
Herbrand-বিচ্ছেদ~$\Sequent !B(t_1) \lor \dots \lor !B(t_n)$-র
!!a{proof} পাই।
% BN-SRC-536 TARGET-CORRECTION-BEGIN
\par\smallskip\noindent\textnormal{\small\textbf{উৎস-সংশোধন (BN-SRC-536):} প্রদর্শিত উপপ্রমাণটির নাম $\pi_1$; উৎসের পরের গদ্যে অনির্ধারিত $\pi_1'$-এর বদলে সেই নাম ব্যবহার করা হয়েছে।}
% BN-SRC-536 TARGET-CORRECTION-END

সমস্যা হলো, $\Sequent \lexists[x][!B(x)]$-র
কর্তনমুক্ত !!{proof}-এও সব \RightR{\lexists} অনুমান
শেষে একসঙ্গে থাকে—এমন ধরে নিতে পারি না। এই $\RightR{\lexists}$ অনুমানগুলির মাঝখানে
কোনো $!B(t_i)$ প্রধান এমন অনুমান থাকতে পারে। এই কারণেই
মধ্যবর্তী সিকোয়েন্ট উপপাদ্য দরকার।

\begin{thm}[মধ্যবর্তী সিকোয়েন্ট উপপাদ্য]\ollabel{thm:midsequent}
$\Gamma \Sequent \Delta$-র এমন \Log{G3c}-!!{proof}~$\pi$
থাকলে, যাতে $\Gamma$ ও $\Delta$-র সব !!{formula}s
অগ্র-পরিমাণসূচক রূপে, এমন !!a{proof}~$\pi'$ আছে যাতে
একটি $\Pi \Sequent \Lambda$ সিকোয়েন্ট থাকে। তার
নীচের সব অনুমান পরিমাণসূচক-সংক্রান্ত, উপরের সব অনুমান
বচনীয়।
\end{thm}

\begin{proof}
  $\pi$-র কোনো পরিমাণসূচক অনুমানের \emph{ক্রম} বলতে তার নীচে
  থাকা বচনীয় অনুমানের সংখ্যা বোঝাই; $\pi$-র ক্রম~$o(\pi)$
  হলো এর সব পরিমাণসূচক অনুমানের ক্রমের যোগফল।
  $o(\pi) = 0$ হলে কোনো পরিমাণসূচক অনুমানের নীচে বচনীয়
  অনুমান নেই, তাই কাজ শেষ। পরিমাণসূচক অনুমান একটিও না
  থাকলে অন্তিম সিকোয়েন্টটিই মধ্যবর্তী সিকোয়েন্ট। অন্যথায়
  সব পরিমাণসূচক অনুমান এক-পূর্বধারণাবিশিষ্ট বলে একটি
  সর্বোচ্চে-থাকা পরিমাণসূচক অনুমান থাকে; তার পূর্বধারণা
  মধ্যবর্তী সিকোয়েন্ট।
  % BN-SRC-537 TARGET-CORRECTION-BEGIN
  \par\smallskip\noindent\textnormal{\small\textbf{উৎস-সংশোধন (BN-SRC-537):} ক্রম শূন্য হলেও পরিমাণসূচক অনুমান একটিও না-থাকতে পারে; সেই ভিত্তিক্ষেত্রে অন্তিম সিকোয়েন্টকেই মধ্যবর্তী ধরা হয়েছে।}
  % BN-SRC-537 TARGET-CORRECTION-END

  $o(\pi) > 0$ হলে ক্রম~$>0$ এমন সর্বনিম্নে-থাকা
  পরিমাণসূচক অনুমান বাছি। তার ক্রম $>0$, তাই নীচে
  অবশ্যই বচনীয় অনুমান আছে। সর্বনিম্ন এমনটি বাছায়
  তার সিদ্ধান্ত অন্য কোনো পরিমাণসূচক অনুমানের পূর্বধারণা
  হতে পারে না: নইলে নীচের সেই পরিমাণসূচক অনুমানেরও
  নীচে বচনীয় অনুমান থাকত। কোন পরিমাণসূচক ও কোন
  পরবর্তী বচনীয় বিধি ব্যবহৃত, তা অনুযায়ী অনেকগুলি ক্ষেত্র।
  অন্তিম সিকোয়েন্টে শুধু অগ্র-পরিমাণসূচক সূত্র থাকায়
  পরিমাণসূচক অনুমানের প্রধান সূত্রটি পরের বচনীয় অনুমানে
  সক্রিয় হতে পারে না। হলে ওই বচনীয় অনুমানের প্রধান
  সূত্রে কোনো বচনীয় সংযোজকের পরিসরের ভিতরে পরিমাণসূচক
  থাকত। উপ-!!{formula} ধর্ম অনুযায়ী সেটি অন্তিম সিকোয়েন্টের
  কোনো !!a{formula}-র উপ-!!{formula} হতে হতো। তাই
  পরিমাণসূচক অনুমানের প্রধান !!{formula} নীচের বচনীয়
  অনুমানের প্রসঙ্গের !!a{element}।

  দুটি উদাহরণ দেখি:

প্রথমে ধরি, $\LeftR{\lif}$-এর ডান পূর্বধারণায়
$\RightR{\lforall}$ শেষ হয়। তাহলে $\pi$-র প্রাসঙ্গিক
উপ-!!{proof}~$\pi_1$-এর শেষাংশ:
\[
\AxiomC{}
\RightLabel{$\pi_2$}
\Deduce$\Gamma' \fCenter \Delta', \lforall[x][!A(x)], !B$
\AxiomC{}
\RightLabel{$\pi_3$}
\Deduce$!C, \Gamma' \fCenter \Delta', !A(c)$
\RightLabel{\RightR{\lforall}}
\UnaryInf$!C, \Gamma' \fCenter \Delta', \lforall[x][!A(x)]$
\RightLabel{\LeftR{\lif}}
\BinaryInf$!B \lif !C, \Gamma' \fCenter \Delta', \lforall[x][!A(x)]$
\DisplayProof
\]
$\pi$ নিয়মিত ধরি। তাই ঈজেনচলক~$c$ $\pi_3$-র বাইরে
ঘটে না; বিশেষত $!B \lif !C$ বা $\pi_2$-তে নেই।
এবার !!{proof}~$\pi_1'$ দেখি:
\[
\AxiomC{}
\RightLabel{$\pi_2'$}
\Deduce$\Gamma' \fCenter \Delta', \lforall[x][!A(x)], !A(c), !B$
\AxiomC{}
\RightLabel{$\pi_3$}
\Deduce$!C, \Gamma' \fCenter \Delta', \lforall[x][!A(x)], !A(c)$
\RightLabel{\LeftR{\lif}}
\BinaryInf$!B \lif !C, \Gamma' \fCenter \Delta', \lforall[x][!A(x)], !A(c)$
\RightLabel{\RightR{\lforall}}
\UnaryInf$!B \lif !C, \Gamma' \fCenter \Delta', \lforall[x][!A(x)], \lforall[x][!A(x)]$
\DisplayProof
\]
এখানে $\pi_2$-র ডান পাশে $!A(c)$ যোগ করে দুর্বলীকরণে
$\pi_2'$ পাওয়া যায়। (এতে নতুন অনুমান যোগ হয় না; বিশেষত
ক্রম পাল্টাতে পারে এমন পরিমাণসূচক অনুমানও নয়।
$\pi_2$-তে ঈজেনচলকের শর্তও ভাঙে না, কারণ $!A(c)$-এর
কোনো ধ্রুবক $\pi_2$-তে ঈজেনচলক হিসেবে ব্যবহৃত নয়।)
\RightR{\lforall}-এর ঈজেনচলক-শর্ত এখনও পূরণ হয়, কারণ
$!B \lif !C$-তে~$c$ নেই।

সংকোচনের গ্রহণযোগ্যতা প্রয়োগে $!B \lif !C, \Gamma'
\fCenter \Delta', \lforall[x][!A(x)]$-র
!!a{proof}~$\pi_1''$ পাই। $\lforall[x][!A(x)]$-এর
জন্য $\RightR{\Contraction}$-এর গ্রহণযোগ্যতার প্রমাণে,
এবং সেখানে ব্যবহৃত $\RightR{\lforall}$-এর উল্টনযোগ্যতার
প্রমাণে, কোনো অনুমানের ক্রম বদলায়নি: বড়জোর অনুমান
বাদ পড়েছে। তাই $\pi_1''$-তে $\pi_1$-এর চেয়ে বেশি
পরিমাণসূচক অনুমান নেই; আর $\pi_1''$-এর প্রতিটির নীচে
বচনীয় অনুমানের সংখ্যা $\pi_1$-এ সংশ্লিষ্ট অনুমানের
নীচের সংখ্যার সমান—শেষ \RightR{\lforall} ব্যতিক্রম।
$\pi$-তে তার নীচে একটি \LeftR{\lif} ছিল, যা এখন
$\pi_1''$-তে তার উপরে সরানো হয়েছে। $\pi$-তে $\pi_1$
বদলে $\pi_1''$ বসাই। নতুন !!{proof}-এর ক্রম কম,
কারণ অন্তত ওই $\RightR{\lforall}$ অনুমানের ক্রম
(অন্তত)~$1$ কমেছে। এবার আরোহের অনুমান প্রয়োগ করি।

পরিমাণসূচক অনুমান \RightR{\lexists} হলে ক্ষেত্রগুলি আরও
সহজ: সংকোচন লাগে না, ঈজেনচলকের শর্তও ভাবতে হয় না।
যেমন, ধরি \RightR{\lexists}-এর পরে $\RightR{\land}$
আসে (এবার বাম পূর্বধারণায়)। প্রাসঙ্গিক উপ-!!{proof} হলো~$\pi_1$:
\[
\AxiomC{}
\RightLabel{$\pi_2$}
\Deduce$\Gamma' \fCenter \Delta', \lexists[x][!A(x)], !A(t), !B$
\RightLabel{\RightR{\lexists}}
\UnaryInf$\Gamma' \fCenter \Delta', \lexists[x][!A(x)], !B$
\AxiomC{}
\RightLabel{$\pi_3$}
\Deduce$\Gamma' \fCenter \Delta', \lexists[x][!A(x)], !C$
\RightLabel{\RightR{\land}}
\BinaryInf$\Gamma' \fCenter \Delta', \lexists[x][!A(x)], !B \land !C$
\DisplayProof
\]
% BN-SRC-538 TARGET-CORRECTION-BEGIN
\par\smallskip\noindent\textnormal{\small\textbf{উৎস-সংশোধন (BN-SRC-538):} এই প্রমাণবৃক্ষের মধ্যবর্তী পূর্বপক্ষে অতিরিক্ত বিস্ময়চিহ্ন ছিল; অন্য তিনটি সিকোয়েন্টের সঙ্গে মিলিয়ে সেটি বাদ দেওয়া হয়েছে।}
% BN-SRC-538 TARGET-CORRECTION-END
$\pi$-তে $\pi_1$-এর বদলে $\pi_1'$ বসাই:
\[
\AxiomC{}
\RightLabel{$\pi_2$}
\Deduce$\Gamma' \fCenter \Delta', \lexists[x][!A(x)], !A(t), !B$
\AxiomC{}
\RightLabel{$\pi_3'$}
\Deduce$\Gamma' \fCenter \Delta', \lexists[x][!A(x)], !A(t), !C$
\RightLabel{\RightR{\land}}
\BinaryInf$\Gamma' \fCenter \Delta', \lexists[x][!A(x)], !A(t), !B \land !C$
\RightLabel{\RightR{\lexists}}
\UnaryInf$\Gamma' \fCenter \Delta', \lexists[x][!A(x)], !B \land !C$
\DisplayProof
\]
এখানে $\pi_3$-র ডান পাশে $!A(t)$ যোগ করে দুর্বলীকরণে
$\pi_3'$ পাই। শুধু $\pi_3$-র প্রত্যেক সিকোয়েন্টের ডান
পাশে $!A(t)$ যোগ হয়; তাই কোনো পরিমাণসূচক অনুমানের
ক্রম বদলায় না। স্থানবিনিময়-করা \RightR{\lexists}
অনুমানটির ক্রম~$1$ কমে, আর $\pi'$-তে বাকি অনুমানগুলির ক্রম
$\pi$-র মতোই থাকে। আরোহের অনুমান প্রযোজ্য।
\end{proof}

\begin{prob}
\olref{thm:midsequent}-এর প্রমাণে \LeftR{\lexists}
যখন \RightR{\lif}-এর পূর্বধারণায় শেষ হয়, এবং
\LeftR{\lforall} যখন \LeftR{\lor}-এর বাম পূর্বধারণায়
শেষ হয়—এই ক্ষেত্রগুলি বর্ণনা করুন।
\end{prob}

\begin{proof}[Herbrand-এর উপপাদ্যের প্রমাণ]
ধরি, $\pi$-র অন্তিম সিকোয়েন্ট $\Sequent \lexists[x][!A(x)]$।
\olref{thm:midsequent} অনুযায়ী $\pi$-কে এমন
!!a{proof}~$\pi'$-তে বদলাই যার একটি মধ্যবর্তী
সিকোয়েন্ট~$S$ আছে। তার রূপ অবশ্যই
\[\Sequent \lexists[x][!A(x)], !A(t_1), \dots, !A(t_n).\]
$S$-এর উপরের সব অনুমান বচনীয় বলে
$\lexists[x][!A(x)]$ $S$-এর উপরের কোনো অনুমানের প্রধান
সূত্র হতে পারে না। এটি পরমাণবিক নয়, তাই স্বতঃসিদ্ধেও
প্রধান হতে পারে না। $!A(t_1)$-এর প্রকৃত উপসূত্রও
$\lexists[x][!A(x)]$ হতে পারে না (মনে রাখি,
$!A(x)$ পরিমাণসূচকবিহীন); তাই $S$-এর উপরেও এটি
সক্রিয় হতে পারে না। ফলে~$S$-তে শেষ হওয়া
উপ-!!{proof} থেকে $\lexists[x][!A(x)]$ বাদ দিলে
$\Sequent !A(t_1), \dots, !A(t_n)$-র সঠিক
!!{proof} পাই। সেখান থেকে $n-1$টি
$\RightR{\lor}$ অনুমানে
$\Sequent !A(t_1) \lor \dots \lor !A(t_n)$-র
!!a{proof} পাওয়া যায়।
\end{proof}

% \begin{prob}
% $\lexists[x][\lforall[y][(!A(x) \lif
% !A(y))]]$-এর একটি Herbrand-বিচ্ছেদ নির্ণয় করুন: $\Sequent
% \lexists[x][\lforall[y][(!A(x) \lif !A(y))]]$-র
% কর্তনমুক্ত !!{proof} খুঁজে এবং \olref{thm:midsequent}-এর
% রূপান্তরে প্রয়োজনে মধ্যবর্তী সিকোয়েন্ট বের করুন।
% \end{prob}

\end{document}
